% Generated draft: requires mathematical and source review.
\documentclass[11pt]{article}
\usepackage[margin=25mm]{geometry}
\usepackage{fontspec}
\setmainfont{Latin Modern Roman}
\usepackage{amsmath,amssymb,mathtools,unicode-math}
\setmathfont{Latin Modern Math}
\usepackage{enumitem,needspace,xurl,hyperref}
\setlength{\parindent}{0pt}
\setlength{\parskip}{5pt}
\setlength{\emergencystretch}{3em}
\newcommand{\N}{\mathbb{N}}
\newcommand{\Z}{\mathbb{Z}}
\newcommand{\Q}{\mathbb{Q}}
\newcommand{\R}{\mathbb{R}}
\newcommand{\C}{\mathbb{C}}
\newcommand{\F}{\mathbb{F}}
\newcommand{\A}{\mathbb{A}}
\newcommand{\PP}{\mathbb{P}}
\newcommand{\E}{\mathbb{E}}
\newcommand{\eps}{\varepsilon}
\newcommand{\dd}{\,\mathrm{d}}
\begin{document}
\section*{UnsolvedMath ID: 700409}
\textbf{Source problem code:} catalogue entry
\subsection{Definitions and mathematical statement}
Let \(k\in\Z_{\geq 1}\), let \(n\in\Z\) be odd with \(n\geq 665\), and let \(F_1,\ldots,F_k\) be nontrivial groups containing no involutions. An involution is an element of order exactly \(2\). Thus, for each \(i\), the group \(F_i\) has an element other than its identity, and \(x^2=1\) implies \(x=1\) for every \(x\in F_i\).

Write
\[
P=\operatorname{Per}_n(F_1,\ldots,F_k)
\]
for the periodic product of exponent parameter \(n\) intended in S. I. Adian's problem. This notation names the source construction; the supplied record does not define that construction. In particular, no exponent assumption on the individual groups \(F_i\) is stated, and no such assumption is added here. A complete definition of this periodic product must be supplied before this draft is a self-contained formulation.

For a group \(G\) and a subset \(X\subseteq G\), let \(\langle X\rangle\) denote the smallest subgroup of \(G\) containing \(X\). The requested assertion, subject to the construction qualification above, is that for every choice of \(k,n,F_1,\ldots,F_k\) satisfying the stated hypotheses,
\[
\forall X\subseteq P,\qquad |X|<k\ \Longrightarrow\ \langle X\rangle\neq P.
\]
Here \(\langle\varnothing\rangle=\{1\}\), so the assertion for \(k=1\) says that \(P\) is nontrivial. The assertion is a lower bound on the number of generators; it does not assert that \(P\) itself is generated by \(k\) elements.

The source additionally states that this assertion, together with the cited work of Adian from 1978, would imply
\[
\forall k\in\Z_{\geq1}\ \exists S\quad
\left\{
\begin{aligned}
&S\text{ is a nontrivial simple group},\\
&\exists Y\subseteq S:\ |Y|=k\text{ and }\langle Y\rangle=S,\\
&\forall X\subseteq S:\ |X|<k\Longrightarrow\langle X\rangle\neq S.
\end{aligned}
\right.
\]
A nontrivial group is simple if its only normal subgroups are the identity subgroup and the whole group. This implication is reported as a source qualification; its supporting hypotheses and argument are not supplied in the record.
\subsection{Short English statement}
Adian asks whether every periodic product with odd exponent parameter \(n\geq665\) of \(k\) nontrivial groups without involutions requires at least \(k\) generators. The source says this would yield simple groups requiring exactly \(k\) generators for every positive integer \(k\). The intended periodic product needs a complete definition from the source literature.
\subsection{Sources}
\begin{itemize}
\item Local UnsolvedMath registry: permanent numeric ID \(700409\), exact source code \texttt{catalogue entry}, title \emph{Kourovka Notebook Problem 7.3}. The supplied registry labels the problem \texttt{open}; this label is supplied metadata, not an independently verified mathematical or literature assessment.
\item \emph{The Kourovka Notebook (Unsolved Problems in Group Theory)}, original problem 7.3, seventh issue (1980). Proposed by S. I. Adian. The supplied background identifies the twenty-first issue edition (2026), \texttt{arXiv:1401.0300v48}. The collection metadata names E. I. Khukhro and V. D. Mazurov as editors.
\item S. I. Adian, \emph{Sov. Math. Doklady}, volume 19 (1978), pages 910--913. The source cites this work as supporting the claimed implication concerning simple groups. No article title or further bibliographic details are supplied.
\end{itemize}
\end{document}
